\documentclass[11pt]{article}
\usepackage[T1]{fontenc}
\usepackage{lmodern}
\usepackage[margin=0.85in]{geometry}
\usepackage{hyperref}
\pagestyle{empty}

\begin{document}

\begin{center}

\vspace*{0.1cm}

{\LARGE\bfseries Recalibrating Tail Risk Forecasts\\[4pt]
under Temporal Dependence}

\vspace{0.4cm}

{\large
Daniel Traian Pele$^{a,b,c,*}$\,,\quad
Vlad Bolov\u{a}neanu$^{a}$\,,\quad
Andrei Theodor Ginavar$^{a,c}$\,,\\[4pt]
Stefan Lessmann$^{a,d}$\,,\quad
Wolfgang Karl H{\"a}rdle$^{a,c,d,e,f}$
}

\vspace{0.25cm}

\begin{flushleft}
\small
$^{a}$\,Bucharest University of Economic Studies (ASE), Bucharest, Romania\\[1pt]
$^{b}$\,Institute for Economic Forecasting, Romanian Academy, Bucharest, Romania\\[1pt]
$^{c}$\,IDA Institute of Digital Assets, Bucharest, Romania\\[1pt]
$^{d}$\,School of Business and Economics, Humboldt University of Berlin, Berlin, Germany\\[1pt]
$^{e}$\,School of Informatics, University of Edinburgh, Edinburgh, UK\\[1pt]
$^{f}$\,Dept.\ of Information Management and Finance, National Yang Ming Chiao Tung University, Hsinchu, Taiwan\\[4pt]
$^{*}$\,Corresponding author: \href{mailto:danpele@ase.ro}{danpele@ase.ro}
\end{flushleft}

\end{center}

\vspace{0.15cm}

\noindent\textbf{Abstract}

\smallskip

\noindent
Extreme-quantile forecasting is constrained by tail sparsity: at the 1\% level, only a small number of exceedances are available for calibration, making flexible post-processing statistically fragile. We propose a low-dimensional alternative: a one-parameter conformal shift applied to the predicted lower quantile. Beyond recalibration, the shift is used as a signed audit statistic, measuring whether the base forecast retains usable tail information or whether the correction effectively replaces it.

Applied to classical volatility models and zero-shot time-series foundation models, the audit reveals a sharp regime separation. In the \emph{signal-preserving} regime, the conformal shift makes a small adjustment to an informative forecast; in the \emph{replacement} regime, the correction dominates the raw tail estimate. This distinction exposes when apparent calibration gains reflect genuine forecast information and when they reflect empirical fallback.

Rolling recalibration substantially improves extreme-tail coverage and outperforms richer recalibration alternatives, while fitted per-asset GARCH models retain a Quantile Score advantage. The results show that local, low-dimensional recalibration can be effective for marginal tail calibration, but they also identify its boundary: scalar shifts do not restore conditional violation independence. Expected Shortfall is therefore treated as a supplementary diagnostic rather than as a target with formal conformal guarantees.

\medskip

\noindent\textbf{Keywords:} Conformal prediction; Value-at-Risk; Time series foundation models; Tail risk forecasting; Temporal dependence; Forecast audit; Quantile Score

\smallskip

\noindent\textbf{JEL Classification:} C14\,|\,C53\,|\,C58\,|\,G17\,|\,G32

\vspace{0.4cm}

\noindent\textbf{Acknowledgements}

\smallskip

\noindent\footnotesize
This paper is supported through the project ``IDA Institute of Digital Assets'', CF166/15.11.2022, CN760046/23.05.2023; the project ``AI for Energy Finance (AI4EFin)'', CF162/15.11.2022, CN760048/23.05.2023, financed under Romania's National Recovery and Resilience Plan, Call no.\ PNRR-III-C9-2022-I8; and the Marie Sk\l{}odowska-Curie Actions under the European Union's Horizon Europe research and innovation program for the Industrial Doctoral Network on Digital Finance, acronym DIGITAL, Project No.\ 101119635.
We acknowledge the support of the project ``MA'AT -- Autonomous Model for Textual Assistance'', SMIS Code 2021+: 330941, funding contract no.\ 390090/11.11.2025, project co-financed by the European Regional Development Fund through the Smart Growth, Digitalisation and Financial Instruments Programme 2021--2027 (POCIDIF).

\end{document}
